% !PRINTED 382
\weberpage{382}
%<<CONT
deren Wurzeln die zehn mit $\xi_1\xi_2$ syzygetischen Wurzeln sind. Bedeutet $\xi_3$ eine von diesen, so ist also

\weq{3}{X(\xi_1,\xi_2,\xi_3) = 0,}

und es folgt daraus der Satz

\gesperrt{1. Es giebt eine rationale Function von drei Variablen $X(\xi_1,\xi_2,\xi_3)$, die verschwindet, wenn $\xi_1\xi_2\xi_3$ irgend ein syzygetisches Tripel von Wurzeln der Doppeltangentengleichung ist, und die nicht verschwindet, wenn $\xi_1\xi_2\xi_3$ ein azygetisches Tripel ist.}

Nach \S.\,100 können durch ein vollständiges Siebener-System dieser Wurzeln

\weq{4}{\xi_1,\ \xi_2,\ \xi_3,\ \xi_4,\ \xi_5,\ \xi_6,\ \xi_7}

alle Wurzeln rational ausgedrückt werden, und zwar in der Weise, dass z.\,B.

\weq{5}{\xi_{12} = \Psi(\xi_1,\xi_2 \mid \xi_3,\xi_4,\xi_5,\xi_6,\xi_7)}

eine Wurzel wird, wo $\Psi$ eine rationale Function bedeutet, die sich nicht ändert, wenn $\xi_1$ und $\xi_2$ vertauscht oder wenn $\xi_3,\xi_4,\xi_5,\xi_6,\xi_7$ beliebig permutirt werden. Wird aber in (5) bei festgehaltener Function $\Psi$ an Stelle von $\xi_1,\xi_2$ ein anderes Paar $\xi_i,\xi_k$ gesetzt, so erhält man eine andere Wurzel $\xi_{ik}$. Die Wurzeln (4) sind auch mit $[1\,8],\dots[7\,8]$ und $\xi_{ik}$ mit $[i\,k]$ zu bezeichnen.

Durch jede Permutation der Wurzeln (1) wird nach (2) das ganze System der Wurzeln $[i\,k]$ eine gewisse Permutation erfahren.

Ersetzt man aber das vollständige Siebener-System (4) durch ein anderes, so ergiebt die Formel (5) eine bestimmte andere Wurzel, und das ganze System der Wurzeln $[i\,k]$ wird einer zweiten Permutation unterworfen. Es ist dann zunächst leicht zu beweisen:

\gesperrt{2. Die Permutationsgruppe $P$ der Wurzeln der Doppeltangentengleichung, die man erhält, wenn man in (4) und (5) an Stelle von $\xi_1,\xi_2,\dots\xi_7$ alle vollständigen Systeme, jedes in jeder beliebigen Ordnung, setzt, ist die Galois'sche Gruppe der Doppeltangentengleichung.}

Um dies nachzuweisen, haben wir Zweierlei zu zeigen:
%CONT>>
